% GENERATED FILE. Edit canonical lessons, then run npm run generate:books.
% canonical-source-hash: fnv1a64:6e7d561458eb282d
% canonical-lessons: IT-C387-requisito, IT-C387-idoneo, IT-C387-ammissione, IT-C387-escludere, IT-C387-condizione

\chapter{Chi può partecipare}
\label{ch:it-chi-puo-partecipare}
\hlchaptermodality{387}

\begin{chapteropening}
\textbf{By the end of this chapter:} \emph{I can say who can take part: the requirements, being eligible, admission, being left out and the conditions.}

Complete the last lesson of chapter 387: I can say who can take part: the requirements, being eligible, admission, being left out and the conditions.
\end{chapteropening}

\section[il requisito]{il requisito — the requirement}

\label{lesson:IT-C387-requisito}



\begin{quote}
Before the new one: say the Italian for to break, then the Italian for the exception.
\end{quote}

\subsection*{You'll want to know: il requisito}
\textbf{il requisito} — \textquotedblleft{}the requirement\textquotedblright{}.

\begin{quote}
Il primo requisito è avere la patente.
\end{quote}

\begin{quote}
\emph{The first requirement is to have a driving licence.}
\end{quote}

Mostly plural: \emph{i requisiti}.

\begin{cousinweb}
\emph{requisito} is a learned borrowing from Latin \emph{requīsītum} \textquotedblleft{}a thing needed, a demand\textquotedblright{} (Lewis \& Short), from \emph{requīrere} \textquotedblleft{}to seek again, to need\textquotedblright{}, built on \emph{quaerere} \textquotedblleft{}to seek\textquotedblright{}.
\end{cousinweb}

\begin{grammarlens}[title={What you've built}]
One more word for telling, explaining and describing at B1.
\end{grammarlens}

\subsection*{Your turn}
\begin{itemize}
\raggedright
  \item \emph{Say it:} \emph{il requisito}
  \item \emph{Say it:} \emph{il requisito}, once more
  \item \emph{Say it:} a sentence of your own with \emph{il requisito}
  \item \emph{Recall:} say the Italian for to break, then the Italian for the exception, then say \emph{il requisito} again
\end{itemize}

\begin{tcolorbox}[breakable,colback=teal!4,colframe=teal!35!black,title={Before you move on}]
What does il requisito mean? (\textquotedblleft{}The requirement\textquotedblright{}.) Say it once more.
\end{tcolorbox}
\section[idoneo]{idoneo — suitable, eligible, fit}

\label{lesson:IT-C387-idoneo}



\begin{quote}
Before the new one: say the Italian for the exception, then the Italian for the requirement.
\end{quote}

\subsection*{You'll want to know: idoneo}
\textbf{idoneo} — \textquotedblleft{}suitable, eligible, fit\textquotedblright{}.

\begin{quote}
Se il medico ti giudica idoneo, puoi fare la gara.
\end{quote}

\begin{quote}
\emph{If the doctor finds you fit, you can take part in the race.}
\end{quote}

Formal: in exams \emph{idoneo} means passed, without a mark.

\begin{cousinweb}
\emph{idoneo} is a learned borrowing from Latin \emph{idōneus} \textquotedblleft{}fit, suitable\textquotedblright{} (Lewis \& Short); the deeper origin of the Latin word is not settled.
\end{cousinweb}

\begin{grammarlens}[title={What you've built}]
One more word for telling, explaining and describing at B1.
\end{grammarlens}

\subsection*{Your turn}
\begin{itemize}
\raggedright
  \item \emph{Say it:} \emph{idoneo}
  \item \emph{Say it:} \emph{idoneo}, once more
  \item \emph{Say it:} a sentence of your own with \emph{idoneo}
  \item \emph{Recall:} say the Italian for the exception, then the Italian for the requirement, then say \emph{idoneo} again
\end{itemize}

\begin{tcolorbox}[breakable,colback=teal!4,colframe=teal!35!black,title={Before you move on}]
What does idoneo mean? (\textquotedblleft{}Suitable, eligible, fit\textquotedblright{}.) Say it once more.
\end{tcolorbox}
\section[l'ammissione]{l'ammissione — the admission, the entry}

\label{lesson:IT-C387-ammissione}



\begin{quote}
Before the new one: say the Italian for the requirement, then the Italian for suitable.
\end{quote}

\subsection*{You'll want to know: l'ammissione}
\textbf{l'ammissione} — \textquotedblleft{}the admission, the entry\textquotedblright{}.

\begin{quote}
L'esame di ammissione è a settembre.
\end{quote}

\begin{quote}
\emph{The entrance exam is in September.}
\end{quote}

Built on \emph{ammettere}. It is also an admission of guilt.

\begin{cousinweb}
\emph{ammissione} is a learned borrowing from Latin \emph{admissiō} \textquotedblleft{}a letting in\textquotedblright{} (Lewis \& Short), from \emph{admittere} \textquotedblleft{}to let in\textquotedblright{}, built on \emph{mittere} \textquotedblleft{}to send\textquotedblright{}.
\end{cousinweb}

\begin{grammarlens}[title={What you've built}]
One more word for telling, explaining and describing at B1.
\end{grammarlens}

\subsection*{Your turn}
\begin{itemize}
\raggedright
  \item \emph{Say it:} \emph{l'ammissione}
  \item \emph{Say it:} \emph{l'ammissione}, once more
  \item \emph{Say it:} a sentence of your own with \emph{l'ammissione}
  \item \emph{Recall:} say the Italian for the requirement, then the Italian for suitable, then say \emph{l'ammissione} again
\end{itemize}

\begin{tcolorbox}[breakable,colback=teal!4,colframe=teal!35!black,title={Before you move on}]
What does l'ammissione mean? (\textquotedblleft{}The admission, the entry\textquotedblright{}.) Say it once more.
\end{tcolorbox}
\section[escludere]{escludere — to exclude, to leave out; to rule out}

\label{lesson:IT-C387-escludere}



\begin{quote}
Before the new one: say the Italian for suitable, then the Italian for the admission.
\end{quote}

\subsection*{You'll want to know: escludere}
\textbf{escludere} — \textquotedblleft{}to exclude, to leave out; to rule out\textquotedblright{}.

\begin{quote}
Se arrivi tardi, ti escludono dalla gara.
\end{quote}

\begin{quote}
\emph{If you arrive late, they will exclude you from the race.}
\end{quote}

Past participle \emph{escluso}. \emph{Lo escludo} — I rule it out.

\begin{cousinweb}
\emph{escludere} is a learned borrowing from Latin \emph{exclūdere} \textquotedblleft{}to shut out\textquotedblright{} (Lewis \& Short), built from \emph{ex} \textquotedblleft{}out\textquotedblright{} and \emph{claudere} \textquotedblleft{}to shut\textquotedblright{}, the verb behind \emph{chiudere}.
\end{cousinweb}

\begin{grammarlens}[title={What you've built}]
One more word for telling, explaining and describing at B1.
\end{grammarlens}

\subsection*{Your turn}
\begin{itemize}
\raggedright
  \item \emph{Say it:} \emph{escludere}
  \item \emph{Say it:} \emph{escludere}, once more
  \item \emph{Say it:} a sentence of your own with \emph{escludere}
  \item \emph{Recall:} say the Italian for suitable, then the Italian for the admission, then say \emph{escludere} again
\end{itemize}

\begin{tcolorbox}[breakable,colback=teal!4,colframe=teal!35!black,title={Before you move on}]
What does escludere mean? (\textquotedblleft{}To exclude, to leave out; to rule out\textquotedblright{}.) Say it once more.
\end{tcolorbox}
\section[la condizione]{la condizione — the condition}

\label{lesson:IT-C387-condizione}



\begin{quote}
Before the new one: say the Italian for the admission, then the Italian for to exclude.
\end{quote}

\subsection*{You'll want to know: la condizione}
\textbf{la condizione} — \textquotedblleft{}the condition\textquotedblright{}.

\begin{quote}
Ti aiuto a una condizione: poi tu aiuti me.
\end{quote}

\begin{quote}
\emph{I will help you on one condition: then you help me.}
\end{quote}

\emph{A condizione che} + subjunctive — on condition that. \emph{In buone condizioni} — in good shape.

\begin{cousinweb}
\emph{condizione} is a learned borrowing from Latin \emph{condiciō}, later spelled \emph{conditiō}, \textquotedblleft{}agreement, terms\textquotedblright{} (EtymDB), from \emph{condīcere} \textquotedblleft{}to agree upon\textquotedblright{}, built on \emph{dīcere} \textquotedblleft{}to say\textquotedblright{}.
\end{cousinweb}

\begin{grammarlens}[title={What you've built}]
One more word for telling, explaining and describing at B1.
\end{grammarlens}

\subsection*{Your turn}
\begin{itemize}
\raggedright
  \item \emph{Say it:} \emph{la condizione}
  \item \emph{Say it:} \emph{la condizione}, once more
  \item \emph{Say it:} a sentence of your own with \emph{la condizione}
  \item \emph{Recall:} say the Italian for the admission, then the Italian for to exclude, then say \emph{la condizione} again
\end{itemize}

\begin{tcolorbox}[breakable,colback=teal!4,colframe=teal!35!black,title={Before you move on}]
What does la condizione mean? (\textquotedblleft{}The condition\textquotedblright{}.) Say it once more.
\end{tcolorbox}
